\problemname{Group Photo}

\illustration{0.3}{pexels-pnw-prod-7625047-small.jpg}{
    Illustration of Sample Input 1.
    Free Pexels License by~PNW~Production on \href{https://www.pexels.com/photo/people-supporting-each-other-7625047/}{Pexels}
}

% optionally define variables/limits for this problem
\newcommand{\maxn}{5 \cdot 10^5}

The members of the No-Weather-too-Extreme Recreational Climbing society
completed their $200$th successful summit today! To commemorate the occasion,
you will take a picture of all the members standing together in one row.

After the photo for the $100$th summit eight years ago turned into a moderate
fiasco, you are determined to get things right this time and make sure that
people are arranged in an aesthetically pleasing way \textit{before} taking the
photo.

As this group is all about climbing mountains, you want the heights of the
climbers in the photo to make the shape of a mountain. More
precisely, the climbers
should be arranged such that their heights are first increasing and then
decreasing (the increasing or decreasing part is allowed to be empty). The
heights of all climbers are pairwise distinct, so for simplicity we will say
that the shortest climber has height $1$, the second-shortest climber has
height $2$, and so on.

The climbers have already positioned themselves in a row in some arbitrary way
that is not necessarily visually pleasing. You will select a subset of the
climbers and rearrange their positions among themselves, with all other climbers
staying exactly where they are. To keep chaos to a minimum, you want the number
of moving climbers to be as small as possible.

What is the size of the smallest subset of climbers such that it is possible
to reorder them so that the sequence of heights becomes first increasing and then
decreasing?

\begin{Input}
    The input consists of:
    \begin{itemize}
        \item One line with an integer $n$ ($1\leq n\leq \maxn$), the number of
			climbers.
        \item One line with $n$ distinct integers $a_1, \ldots, a_n$ ($1 \leq a_i \leq n$ for each $i$), where $a_i$ is the height of the $i$th climber in the current arrangement.
        % $a$ ($1 \leq a \leq n$), the
		% 	heights of the climbers as they are currently arranged in the row.
    \end{itemize}
\end{Input}

\begin{Output}
    Output the minimum number of climbers who need to move.
\end{Output}

\nextsample

Swapping the first and the last climber is the optimal way
to make the sequence of heights first increasing and then decreasing.
